% answers_short.tex —— 1977 年·福建卷·文科 · 简短答案速查
% 由 tools/make_short_answers.py 自 answers.tex 的 \gksolhead/\gkanswer 提取，
% 并与 1977/1977_数学真题_福建文_答案要点.md 逐题比对后生成，勿手改（改源 MD 后重跑）。

\begin{gknotice}
\textbf{使用说明}：本节只给最终答案便于快速核对，完整推导见解析版。答案为编写者独立推导，\textbf{非官方标准答案}。
\end{gknotice}

\gkpart{大题}
\begin{itemize}[leftmargin=2.6em]
  \item[\textbf{1.}] (1) $7$；(2) $-\dfrac12$；(3) $\lvert2x-3\rvert$；(4) $AM=2\ \mathrm{cm}$；(5) $x=3000$；(6) $-1$；(7) 极小值为 $-5$；(8) $-\dfrac34$；(9) $a_n=-\dfrac29\left(-\dfrac13\right)^{n-1}$
  \item[\textbf{2.}] (1) $(1,2)$；(2) 见解析版；(3) $x=12$；(4) $-2<x<3$；(5) $\dfrac{\sqrt{15}+\sqrt6}{3}$；(6) $216\pi\ \mathrm g\approx678\ \mathrm g$（取 $\pi=3.14$）
  \item[\textbf{3.}] (1) 按十月产量 $1000$ 台、十一月至十二月按相同月增长率递增的条件，可得每月增长率 $10\%$；(2) 源题面未给出一月至十月各月产量或其变化规律，无法由当前信息唯一确定原计划年产量，待补充条件或权威版本
  \item[\textbf{4.}] 第一象限交点为 $(3,3\sqrt3)$；抛物线切线为 $y-3\sqrt3=\dfrac{\sqrt3}{2}(x-3)$，圆切线为 $x+\sqrt3y-12=0$。源题面使用单数“切线方程”，未说明求哪一条，待核。
  \item[\textbf{5.}] 按源式可列 $\dfrac2\alpha-\tan\alpha=1$；若取其锐角数值根，则 $\alpha\approx51.1361^\circ$，$m\approx11.0987$，$r\approx5.7786$。由于源式含 $24\alpha$ 且无答案页，暂作按字面计算的候选值，待权威版本核对。
  \item[\textbf{6.}] 按普通平面角定义，$AB=AC$ 且 $\angle BAC=60^\circ$ 使 $\triangle ABC$ 为等边三角形；此时从 $B$、$C$ 按题给 $120^\circ$、$135^\circ$ 作射线不能得到图示同一交点 $P$，当前源面条件不相容，无法给出可靠唯一距离，待核。
\end{itemize}
